\documentclass[12pt, a4paper]{article} % Classe: article, report o book [8]

% --- Lingua e Codifica (Supporto Italiano) ---
\usepackage[utf8]{inputenc}  % Codifica input
\usepackage[T1]{fontenc}      % Codifica font
\usepackage[italian]{babel}   % Lingua italiana [2]

% --- Impaginazione e Margini ---
\usepackage[margin=2.5cm]{geometry} % Imposta margini (2.5cm su tutti i lati)

% --- Pacchetti Scientifici/Matematici ---
\usepackage{amsmath, amssymb, amsfonts, amsthm} % Matematica avanzata [2]
\usepackage{graphicx} % Per inserire immagini
\usepackage{hyperref} % Per collegamenti ipertestuali
\usepackage{booktabs} % Per tabelle professionali

% --- Formattazione Testo e Altro ---
\usepackage[document]{ragged2e}
\usepackage{xcolor}    % Per usare i colori 

% Titolo e intestazione
\title{Marco Ricci}
\date{30-12-2025}


\begin{document}

\maketitle

\section{Ruolo} 
Marco Ricci è direttore di produzione all'interno dell'azienda e si occupa di :
\begin{itemize}
    \item Pianificazione della produzione.
    \item Coordinamento tra i reparti produttivi.
    \item Gestione delle risorse umane e materiali.
    \item Ottimizzazione dei processi produttivi.
    \item Monitoraggio della qualità del prodotto finito.
\end{itemize}

\subsection{Utilizzo attuale di ACESS}
\begin{itemize}
    \item \textbf{Scopo principale:} Verifica del valore monetario della commessa o dell'RA (Richiesta Assistenza).
    \item \textbf{Utilizzo secondario:} Controllo note tramite anteprima di stampa per l'Ufficio Tecnico (UT) quando mancano nel contratto principale.
\end{itemize}

\subsection{Proposte}
\begin{itemize}
    \item \textbf{Obiettivo:} Creare un dialogo univoco tra Ufficio Tecnico e Commerciale mediante codici univoci utili ed utilizzabili sia dal commerciale che dal tecnico( codici $\mathdollar$).
    \item \textbf{Codici macro-modelli:} Sviluppo di macro-distinte "vuote" ma valorizzate economicamente in tempo reale, utilizzabili dai commerciali per le offerte e successivamente riempite dai tecnici.
    \item \textbf{Gestione contratto impianti:} Creazione di capitoli specifici per le macchine e per le parti di impianto, con possibilità di aggiungere note tecniche dettagliate.
    \item \textbf{Possibilità di notifiche:} Possibilità di gestire le notifiche quando un contratto passa di ufficio.
\end{itemize}

\subsection{Flusso di gestione dei contratti (Workflow)}
\begin{itemize}
    \item \textbf{Ricezione:} Il contratto arriva a Stefania per la registrazione e l'inserimento su Fusion.
    \item \textbf{Smistamento:}
    \begin{itemize}
        \item \textbf{Ricambistica semplice:} Inviata direttamente al reparto dedicato.
        \item \textbf{Commesse/Ampliamenti impegnative:} Gestite dal direttore di produzione.
    \end{itemize}
    \item \textbf{Pianificazione:}
    \begin{itemize}
        \item Utilizzo di un file Excel secondario condiviso con i commerciali per definire date di consegna realistiche prima dell'inserimento ufficiale.
        \item Inserimento in "Stato Impianti" solo dopo la definizione delle date.
        \item Assegnazione al reparto Project Manager.
        \item Torna in ufficio produzione che verifica sia stata inserita in "Stato Impianti" correttamente.
        \item Viene mandato all'ufficio tecnico per la creazione del progetto.
    \end{itemize}
\end{itemize}

\subsection{Criticità e miglioramenti necessari}
\begin{itemize}
    \item \textbf{Qualità del contratto:} 
    \begin{itemize}
        \item Evitare copia-incolla errati da vecchi contratti.
        \item \textbf{Struttura logica:} Le informazioni devono seguire un flusso funzionale (non dati sparsi).
        \item \textbf{Centralizzazione dati:} Unica sezione per date, tensioni, voltaggi e specifiche tecniche.
    \end{itemize}
    \item \textbf{Ricambistica:} Separazione netta dei flussi tra commesse complesse e ricambistica pura, in previsione di un ufficio ricambi dedicato futuro.
\end{itemize}

\subsection{Specifiche tecniche e normative}
\begin{itemize}
    \item \textbf{Sezione dedicata nel contratto:} Necessità di campi specifici per:
    \begin{itemize}
        \item Normative particolari (es. N1090, UL/CSA, UNI 5060).
        \item Capitolati speciali.
        \item Dettagli logistici (Montaggio, trasporto/Incoterms, FAT - Factory Acceptance Test).
    \end{itemize}
    \item \textbf{Gestione "Memoria" clienti:} Integrazione automatica delle specifiche ricorrenti per clienti abituali (es. verniciature o etichettature personalizzate) per evitare ricerche manuali su commesse vecchie.
\end{itemize}

\end{document}